% QMB_n01_grundlage.tex — Kapitel 1: Die geometrische Grundlage
% Quellen: Dok. 230, 307, 314, 365
\chapter{Die geometrische Grundlage}
\label{ch:grundlage}

Die Fundamentalrelation der FFGFT lautet
\begin{equation}
  \tilde T \cdot m = 1,
  \label{eq:tm}
\end{equation}
wobei $\tilde T$ eine kompakte Zeitkoordinate und $m$ die Masse einer
Mode ist.\status{SETZUNG} Diese Relation ist keine abgeleitete
Gleichung, sondern die geometrische Setzung, aus der alles Weitere
folgt.

\section{Der Träger: $T^4/\mathbb{Z}_3$}

Die FFGFT rechnet auf einem kompakten Torus $T^4 = (S^1)^4$ mit einer
$\mathbb{Z}_3$-Gittersymmetrie. Dieser Träger hat zwei wesentliche
Eigenschaften.

\textbf{Kompaktheit.} Jede Koordinate ist periodisch. Das hat eine
unmittelbare Konsequenz: Die zulässigen Moden sind diskret — nicht weil
die Natur diskret postuliert wird, sondern weil Periodizität
diskrete Fouriermoden erzwingt. Die Fouriermoden des Trägers liefern
genau die Teilchenspektren des Standardmodells.\status{K}

\textbf{$\mathbb{Z}_3$-Symmetrie.} Die Quotientierung durch
$\mathbb{Z}_3$ erzeugt neun Fixpunkte und eine Trialitätssymmetrie mit
$|\mathrm{Aut}(D_4)| = 1152 = 2^7\cdot3^2$. Diese Symmetrie wird im
Spektraloperator sichtbar und steuert, welche Moden stabil sind.\status{B}

\section{Der Strukturparameter $\xi$}

Der einzige freie Parameter der Theorie ist
\begin{equation}
  \xi = \frac{4}{30000} = \frac{1}{7500} \approx 1{,}333\times10^{-4}.
\end{equation}
Er ist keine Kalibrierungsgröße, sondern folgt aus den
Galois-Gruppenordnungen des Trägers: Speziell gilt
$\xi = C_2(SU(3)_{\mathrm{fund}})/N_{\mathrm{Fourier}}= (4/3)/10^4$.\status{K}
(Dok.~338, Dok.~353)

Alle Massen, Kopplungen und Korrekturen der FFGFT lassen sich aus $\xi$
ableiten; SI-Ankerwerte (Elektronenmasse, $H_0$) dienen nur der
Umrechnung in SI-Einheiten, nicht der Ableitung.

\section{Die Relation $\tilde T\cdot m = 1$ und ihre Konsequenzen}

Gleichung~\eqref{eq:tm} hat drei unmittelbare Folgen.

\textbf{Modengebundene Zeitskala.} Jede Mode mit Masse $m_i$ trägt
eine eigene kompakte Zeitskala $\tilde T_i = 1/m_i$.\status{B} Die Zeit
ist kein äußerer Parameter, sondern eine modespezifische Eigenschaft des
Trägers. Pauli-Theorem und der fehlende Zeitoperator der Standard-QM
sind in dieser Lesart kein Problem: Auf einem kompakten Träger gilt das
Pauli-Argument nicht (Dok.~334).

\textbf{Massenarithmetik aus Geometrie.} Das Verhältnis
$(m_\mu/m_e)^2 = 43200$ folgt sowohl geometrisch aus den
Wicklungszahlen als auch aus den Galois-Gruppenordnungen
$|GF(9)^*|^2\cdot 5^2\cdot|GF(27)|$.\status{K} (Dok.~338)

\textbf{Fraktale Korrektur.} Der Übergang vom kompakten Träger zur
SI-Skala bringt eine rekursive Korrektur
$K_{\mathrm{frak}} = 1 - 100\xi \approx 0{,}9867$, die aus der
Typ-III-Pullback-Struktur des Trägers folgt.\status{B} (Dok.~333)

\begin{laierspur}
Stell dir vor, Physik passiert nicht auf einer unendlich langen
geraden Linie, sondern auf einer Trommel — einer Oberfläche, die
in sich geschlossen ist. Auf einer Trommel sind nur bestimmte
Schwingungen möglich: die, die nach einem Umlauf wieder dort
ankommen, wo sie angefangen haben. Genau so funktioniert $T^4$:
Nur Moden, die auf dem Torus „passen", existieren — und das sind
exakt die Teilchen, die wir messen.
\end{laierspur}
\quelle{Dok.~230, 307, 314, 338, 353, 365}
